% ---
% school: "college"
% level: "8eme"
% section: null
% chapter: "العمليات والحساب على الأعداد الكسرية"
% source: "Axiomaths"
% title: "الاختزال وتساوي الأعداد الكسرية"
% country: "Tunisie"
% year: 2026
% professor: "Axiomaths"
% difficulty: "Facile"
% points: null
% tags: ["fractions", "simplification", "egalite"]
% ---
% Exercice 2/15 de la série 8ème « العمليات والحساب على الأعداد الكسرية » (statement + \axsol + solution). Compiler avec XeLaTeX + axiomath.sty (+ axiomaths-logo.png).
\documentclass[12pt,a4paper]{article}
\usepackage[level=8b,color=15803D]{axiomath}
\axsetup{title={العمليات والحساب على الأعداد الكسرية}, corrige}

\begin{document}

\begin{corr}[title={الاختزال وتساوي الأعداد الكسرية}, stars=1, source={Axiomaths}]
\begin{enumerate}
  \item اختزل كل عدد كسري إلى كسر غير قابل للاختزال :
  \[\dfrac{42}{56}\ ;\qquad \dfrac{-90}{150}\ ;\qquad \dfrac{-27}{-63}\ ;\qquad \dfrac{35}{-84}\ ;\qquad \dfrac{-132}{-198}\]
  \item أتمم بما يناسب :
  \[\dfrac{3}{7}=\dfrac{\ldots}{35}=\dfrac{-12}{\ldots}\]
  \[\dfrac{-5}{8}=\dfrac{\ldots}{24}=\dfrac{15}{\ldots}\]
  \item هل العددان $\dfrac{-8}{12}$ و $\dfrac{14}{-21}$ متساويان ؟ علّل.
\end{enumerate}
\axsol
\begin{enumerate}[itemsep=14pt]
  \item\leavevmode
  \begin{prop}[قاعدة الاختزال]
  نقسم البسط والمقام على نفس العدد غير المعدوم : \ib{$\dfrac{a\div k}{b\div k}=\dfrac{a}{b}$}
  \end{prop}
  \begin{prop}[قاعدة الإشارة]
  \ib{$\dfrac{-a}{-b}=\dfrac{a}{b}$} \qquad و \qquad \ib{$\dfrac{a}{-b}=\dfrac{-a}{b}=-\dfrac{a}{b}$}
  \end{prop}
  نقسم في كل مرة على القاسم المشترك الأكبر للبسط والمقام :
  \[\dfrac{42}{56}=\dfrac{42\div14}{56\div14}=\dfrac{3}{4}\]
  \[\dfrac{-90}{150}=\dfrac{-90\div30}{150\div30}=\dfrac{-3}{5}=-\dfrac{3}{5}\]
  \[\dfrac{-27}{-63}=\dfrac{27}{63}=\dfrac{27\div9}{63\div9}=\dfrac{3}{7}\]
  \[\dfrac{35}{-84}=-\dfrac{35}{84}=-\dfrac{35\div7}{84\div7}=-\dfrac{5}{12}\]
  \[\dfrac{-132}{-198}=\dfrac{132}{198}=\dfrac{132\div66}{198\div66}=\dfrac{2}{3}\]
  \item نضرب البسط والمقام في نفس العدد غير المعدوم.

  للانتقال من المقام $7$ إلى المقام $35$ نضرب في $5$ ، وللانتقال من البسط $3$ إلى البسط $-12$ نضرب في $-4$ :
  \[\dfrac{3}{7}=\dfrac{3\times5}{7\times5}=\dfrac{15}{35}\]
  \[\dfrac{3}{7}=\dfrac{3\times(-4)}{7\times(-4)}=\dfrac{-12}{-28}\]
  وبنفس الطريقة، من $8$ إلى $24$ نضرب في $3$ ، ومن $-5$ إلى $15$ نضرب في $-3$ :
  \[\dfrac{-5}{8}=\dfrac{-5\times3}{8\times3}=\dfrac{-15}{24}\]
  \[\dfrac{-5}{8}=\dfrac{-5\times(-3)}{8\times(-3)}=\dfrac{15}{-24}\]
  \item نختزل كلا من العددين :
  \[\dfrac{-8}{12}=\dfrac{-8\div4}{12\div4}=\dfrac{-2}{3}\]
  \[\dfrac{14}{-21}=\dfrac{14\div7}{-21\div7}=\dfrac{2}{-3}=-\dfrac{2}{3}\]
  إذن العددان متساويان : \ib{$\dfrac{-8}{12}=\dfrac{14}{-21}=-\dfrac{2}{3}$}.
\end{enumerate}
\end{corr}

\end{document}
